\section{Artifact Summary}

The final benchmark is generated by \texttt{run\_full\_scale\_scale\_law\_exception\_suite.py}. The script streams \texttt{results/full\_scale/condition\_metrics.csv} and updates aggregate summaries online. It does not load the full condition table into memory after generation. This matters because the final condition CSV is large enough to audit but small enough to keep under the repository file-size limit.

The run produced 544,320 compact rows. Each compact row represents 262,144 evaluations and 20,971,520 planning-tick decisions. The represented totals are 142,690,222,080 evaluations and 11,415,217,766,400 planning-tick decisions. The row-count validation file records that the expected row count matches the generated row count.

\section{Why Compact Rows Are Used}

The benchmark is intended to stress the claim, not to simulate every individual episode as a separate line. A fully expanded table with every represented seed, scene, schedule, episode, and planning tick would be unnecessarily large. Compact rows make the condition structure explicit while preserving represented scale through the row weight.

The important property is that every aggregate is computed using the same weight. A policy is not helped by appearing in more rows, and a regime is not helped by a different sampling multiplier. The grid is balanced across task family, regime, scale, data, diversity, contact, latency, safety, and policy.

\section{Determinism}

The suite uses deterministic hashed jitter rather than pseudorandom state. The jitter prevents ties and makes neighboring conditions slightly different, but every value is reproducible from the condition keys. This keeps the benchmark auditable and avoids hidden random seeds.

The deterministic design also makes diffs meaningful. If a formula changes, downstream summary tables change reproducibly. If the script is rerun without code changes, the result should be stable.

\section{Task Factor Rationale}

The task families are chosen to cover several robotics failure surfaces. Pick-and-place is relatively representation-heavy. Peg insertion and tool use are contact-heavy. Cable and deformable manipulation combine contact and morphology. Dexterous in-hand rotation stresses contact, morphology, and action precision. Legged locomotion stresses morphology and timing. Mobile manipulation stresses timing and mixed control. Human-shared workspace action stresses safety.

The task factors are not meant to be exact measurements. They are scenario parameters that make the benchmark cover distinct regimes. The paper therefore reports the factors as part of the benchmark, not as claims about real robot distributions.

\section{Regime Factor Rationale}

The regime axis is the main axis of the paper. The representation-smooth regime is included so the benchmark does not trivially punish scale. Contact-rich, safety-constrained, latency-sensitive, morphology-shifted, and mixed-shift regimes are included because they represent different ways that the bottleneck can move away from model size.

The mixed physical shift regime is deliberately hard. It tests whether a method can handle several bottlenecks at once. Scale-only utility falls near zero in this regime, while the router remains high because it has partial coverage for each bottleneck.

\section{Model Scale Axis}

The model-scale axis uses abstract levels from 0.3B to 30B. These levels should not be read as a claim about any specific architecture. They are used to create a monotonic representation-capacity axis.

The important result is that the scale-only policy improves with model size but does not improve enough. Success rises from 0.496 to 0.519. Utility rises from 0.121 to 0.145. At the same time, scale saturation rises from 0.266 to 0.490. This is the benchmark signature of a scale-law exception: scale helps, but the physical residual grows as the representation bottleneck is reduced.

\section{Data Scale Axis}

The data-scale axis tests whether additional demonstrations or training data can replace physical regime modeling. The data-scaled policy improves over scale-only, reaching aggregate utility 0.159 rather than 0.133. That improvement is real but small relative to the router's utility 0.888.

The conclusion is not that data is unimportant. The conclusion is that data must cover the right physical mechanism. More visual demonstrations cannot fully substitute for force information in discontinuous contact, and more teleoperation data cannot fully substitute for a safety constraint if the action interface cannot enforce it.

\section{Embodiment Diversity Axis}

Embodiment diversity is treated separately from data scale because a dataset can be large but narrow. The embodiment-diverse scaling policy reaches utility 0.294, improving substantially over data-scaled and scale-only baselines. This is a positive result for embodiment scaling.

The result is still conditional. Diversity helps most when the bottleneck is morphology transfer. It helps less when the active bottleneck is latency, safety, or contact. A broad body dataset is therefore not a universal answer; it is a targeted answer to one class of physical shift.

\section{Contact Severity Axis}

Contact severity has three levels: low contact, mode switches, and discontinuous contact. The failure metric is contact failure, not just task failure. This allows the benchmark to distinguish a policy that understands the scene from a policy that handles the force transition.

The contact-instrumented policy reduces contact failure by directly addressing this axis. The router reduces it as well because it can decide when contact is the active bottleneck. The scale-only policy remains exposed because its main gain is representation capacity.

\section{Latency Axis}

Latency has three levels: 20 ms, 80 ms, and 240 ms. The metric is latency miss, which measures temporal mismatch between observation and action. This is different from slow inference in isolation. A slow action can still be safe if the state changes slowly; a fast action can still be wrong if the observation is stale.

The latency/safety guarded policy reduces latency miss. The router also reduces it because it routes conditions with high timing pressure toward a guarded mode. The scale-only policy does not directly address the timing mismatch.

\section{Safety Axis}

Safety tightness is separated from success. This is important because aggregate success can hide unacceptable risk. A policy that completes the task while violating a tight margin should not be treated as equivalent to a safe policy.

The safety metric penalizes overconfidence under physical bottlenecks. Scale-only policies can become more confident as representation improves, but if they lack a safety mechanism, that confidence can increase violation exposure in tight regimes.

\section{Policy Ranking}

The policy ranking follows bottleneck coverage. The oracle has complete coverage and reaches the utility ceiling. The regime-aware router has broad but imperfect coverage and is the best non-oracle policy. The latency/safety guarded policy is strong because safety penalties are costly. The contact-instrumented policy is strong in contact regimes but less general. Embodiment-diverse scaling helps morphology but not all physical axes. Data-scaled and scale-only policies improve representation but leave physical residuals.

This ranking is the core positive finding. The best method is not simply bigger. It is scale plus a regime-conditioned physical action interface.

\section{Why Scale-Only Is Kept Fair}

The benchmark does not make scale-only fail by definition. In representation-smooth settings, scale-only utility is 0.312. Its success improves with model size. Its aggregate utility is positive. These choices prevent the paper from becoming an anti-scaling caricature.

The exception is visible because scale-only is not enough in hard regimes. Mixed physical shift drives utility to 0.001, not because scale is useless, but because several non-representational bottlenecks are active.

\section{Why The Router Is Not An Oracle}

The regime-aware router has residual exception severity 0.095. It does not know the bottleneck perfectly. It only has partial coverage for contact, latency, safety, and morphology. This keeps the router below the oracle and makes the gap meaningful.

The oracle is useful because it shows the upper bound. If a learned method closed the router-oracle gap, it would need better bottleneck inference, better intervention selection, or better physical instrumentation.

\section{Scale Saturation Interpretation}

Scale saturation measures how strongly further model scaling is blocked by unresolved physical bottlenecks. The benchmark shows a counterintuitive but important pattern: as model scale increases, scale-only success improves, but scale saturation also increases. This happens because representation error is no longer the dominant residual.

This is a common pattern in physical systems. Once a model understands the scene well enough, remaining errors expose the action physics. The curve then becomes shallow unless the policy receives a new physical channel.

\section{Exception Severity Interpretation}

Exception severity is not a failure label. It is a measure of residual physical bottleneck after scale-focused improvements. A high exception severity means that the condition should not be summarized by model size alone.

The scale-only policy has exception severity 0.543. The data-scaled policy reduces it to 0.476. The embodiment-diverse policy reduces it to 0.390. Physical interventions reduce it further. The router reduces it to 0.095.

\section{Intervention Recovery Interpretation}

Intervention recovery measures how much utility is recovered by targeting the physical bottleneck. It is high for the oracle, 0.664, and high for the router, 0.526. It is low for scale-only, 0.021. This separation is useful because it distinguishes policies that improve representation from policies that actually address the exception.

\section{Relation To Contact-Rich Manipulation}

Contact-rich manipulation is a canonical exception family. The robot must know not only where an object is but how force, friction, and contact mode will evolve. Scale can improve the semantic prior, but it cannot remove the need for state about contact.

The benchmark therefore separates contact failure from representation error. A policy can have low representation error and still high contact failure. That is exactly the case where a scale-only curve overstates readiness.

\section{Relation To Tactile Sensing}

Tactile sensing is one way to reduce contact uncertainty. The benchmark abstracts this as contact instrumentation rather than modeling a particular sensor. The point is not that tactile sensing is always necessary. The point is that a contact-rich scaling curve should state whether the policy has a physical channel capable of detecting contact state.

\section{Relation To Latency}

Latency creates a temporal exception. If sensing and action are misaligned, a larger model can generate a better answer to the wrong state. The latency/safety guarded policy models compensation. The router models the decision to use such compensation only when latency pressure is active.

\section{Relation To Safety Constraints}

Safety constraints change the objective. A policy can be successful under a loose metric and unacceptable under a tight safety envelope. The benchmark therefore reports safety violation separately and gives it a strong utility penalty.

This is not a claim that the deterministic suite certifies safety. It is a claim that safety should appear as a regime-conditioned scaling slice.

\section{Relation To Morphology Transfer}

Morphology transfer is where embodiment diversity should help most. The benchmark supports that view: embodiment-diverse scaling improves over scale-only. The limitation is that morphology diversity does not automatically solve non-morphology bottlenecks.

This matters for generalist robot training. A dataset can cover many bodies but still miss a timing or contact phenomenon. A useful scaling law should identify which physical variations were actually covered.

\section{Relation To Sim-To-Real}

Sim-to-real transfer is often discussed as a visual or dynamics gap. The benchmark suggests a more decomposed view. The gap may be contact, latency, safety, morphology, or a mixture. A single sim-to-real score can hide which gap is active.

The regime map can therefore be used as an audit tool. If a sim-to-real method improves the aggregate score, the next question is which physical slice improved.

\section{Relation To World Models}

World models can predict hidden state and improve long-horizon planning. They are compatible with the paper's claim. A world model helps most when it predicts the active physical bottleneck in a form that the action policy can use.

The benchmark does not include a separate world-model policy because the main axis is reporting, not architecture. A future extension could compare world-model bottleneck inference against the router and oracle.

\section{Relation To Vision-Language-Action Models}

Vision-language-action models are likely to benefit from scale. They can improve semantic grounding, object recognition, and instruction following. The exception is that many robot failures occur after the instruction is understood.

The benchmark's scale-only policy should be read as a strong VLA-style baseline without explicit physical bottleneck routing. Its positive smooth-regime performance acknowledges the value of VLA scale. Its hard-regime saturation motivates adding physical regime interfaces.

\section{Hostile Interpretation: This Is Just More Features}

A hostile reviewer might say that the router wins because it has more features. That critique is partly fair. The benchmark intentionally gives physical interventions to some policies. The question is whether those interventions are necessary to explain the failure modes. The decomposed metrics show that they are: contact instrumentation reduces contact failure, guarded control reduces latency and safety errors, and embodiment diversity reduces morphology loss.

The router is not rewarded for being complicated in the abstract. It is rewarded when the active bottleneck matches an available physical intervention.

\section{Hostile Interpretation: The Oracle Is Too Strong}

Correct. The oracle is too strong by design. It is an upper bound. It should not be compared to deployable policies as a practical method. Its role is to show how much performance remains if bottleneck inference were perfect.

The more important comparison is between scale-only and the regime-aware router. That comparison does not require the oracle.

\section{Hostile Interpretation: The Benchmark Is Synthetic}

Correct. The benchmark is synthetic and deterministic. It is not a substitute for hardware data. The paper's claim is therefore limited to a benchmark and reporting discipline.

The synthetic benchmark is still useful because it makes the hidden assumption explicit. If a real benchmark disagrees, the disagreement will be informative: either physical bottlenecks are less important than modeled here, or current real metrics are not exposing them.

\section{Falsification Criteria}

The claim would be weakened if scale-only policies matched the router across physical regimes while keeping contact failure, latency miss, safety violation, and morphology loss low. It would also be weakened if real robot logs showed smooth monotonic scale curves after stratifying by contact, latency, safety, and morphology.

The claim would be strengthened if real robot benchmarks showed the same pattern: scale helps in smooth regimes, physical interventions dominate in hard regimes, and route-conditioned policies beat aggregate scale-only curves.

\section{Minimum Real-Robot Follow-Up}

A minimum hardware follow-up would use one manipulation platform and three tasks: one representation-smooth task, one contact-rich task, and one safety- or latency-constrained task. The study would train or evaluate scale-only and physical-intervention policies at multiple data sizes. It would report success, contact failure, latency miss, safety violation, and morphology or calibration loss separately.

The key requirement is not a large number of robots. The key requirement is regime stratification.

\section{Artifact Integrity Checklist}

The final artifact should pass these checks. The PDF should be at least 25 pages. The final PDF should be exported to \texttt{C:/Users/wangz/Downloads/57.pdf} only after finalization. Local \texttt{paper/main.pdf} should be removed. The generated validation JSON should match the represented counts in the paper. Current status docs should mark v3 final status, with v2 preserved only as a negative control. The final hash, file size, page count, and visual QA status should be recorded.

\section{Reproducibility Checklist}

Run \texttt{python run\_full\_scale\_scale\_law\_exception\_suite.py} from the repository root. Then run the final build script after the manuscript has been finalized. The build script should compile from \texttt{paper/}, enforce the page threshold, copy the final artifact to Downloads, record build status, and remove the local PDF.

\section{Why The Result Is Positive}

The result is positive because it does not merely say scale fails. It identifies where scale helps and what must be added when scale saturates. The router's success shows that physical bottleneck awareness can be combined with scaling. The oracle gap shows that better learned bottleneck inference is a meaningful research target.

\section{Why The Result Is Narrow}

The result is narrow because it is not a universal theorem. It is a deterministic benchmark with explicit assumptions. The correct use is to motivate regime-conditioned reporting and future real-robot tests.

\section{Claim Guardrails}

Use the phrase regime-conditioned scaling, not anti-scaling. Use benchmark evidence, not hardware proof. Say that scale-only is positive but bottleneck-limited. Say that the router is best non-oracle in this deterministic suite. Do not claim deployment-ready safety. Do not claim that the exact numeric utilities transfer to real robots.

\section{Reviewer-Ready Summary}

The paper's answer to a skeptical reviewer is concise. V2 showed corpus ratios were not enough. V3 adds a full-scale deterministic benchmark. Scale-only improves with model size but saturates under physical bottlenecks. Physical interventions recover specific failures. The regime-aware router is the best non-oracle policy. The benchmark is synthetic, so the claim is a reporting discipline and research direction, not a hardware certification.

\section{Metric Formula Audit}

The implementation computes physical pressure terms before policy metrics. Contact pressure combines task contact need, regime contact load, contact severity, and base task difficulty. Latency pressure combines task latency sensitivity, regime latency load, latency level, and difficulty. Safety pressure combines task safety pressure, regime safety load, safety tightness, and difficulty. Morphology pressure combines task morphology need, regime morphology load, and embodiment diversity.

This ordering matters. Scale is not allowed to erase the physical pressure term directly. Scale improves representation capacity, and representation capacity can reduce some secondary errors. Physical interventions reduce the relevant physical residuals more directly. That separation is what makes the exception measurable.

\section{Representation Capacity Audit}

Representation capacity is a function of model scale, data scale, embodiment diversity, representation need, and policy focus. The scale-only policy receives the strongest model-scale coefficient. The data-scaled policy receives the strongest data coefficient. The embodiment-diverse policy receives the strongest diversity coefficient.

This is deliberately fair to scaling baselines. If scale-only still saturates, the reason is not that the benchmark denied it representation improvement. The reason is that remaining utility is controlled by physical residuals.

\section{Physical Coverage Audit}

Physical coverage is policy-specific. Contact instrumentation covers contact pressure. Latency/safety guarding covers timing and risk pressure. Embodiment diversity covers morphology pressure. The router has moderate coverage on all physical axes and an additional router bonus. The oracle has high coverage on every physical axis.

This design supports mechanism-level interpretation. If a contact-instrumented policy improves contact failure but not latency miss, that is expected. If a router improves several errors at once, it is because it has broad partial coverage.

\section{Overconfidence Term}

The overconfidence term is included to model a common deployment failure. A strong model can act decisively even when the active bottleneck is not represented. This term is small, but it raises safety and contact risk for scale-focused policies in high-bottleneck regimes.

The point is not that large models are inherently unsafe. The point is that confidence should be calibrated to the physical regime. A model that knows what the object is still needs to know whether the contact, timing, and safety state make the action valid.

\section{Why Utility Is Cost-Aware}

The benchmark reports success, but success is not enough. A policy that completes a task while violating a safety margin or wasting repeated physical searches should not be treated as equivalent to a safer and more efficient policy.

Utility therefore rewards success, intervention recovery, and data efficiency while penalizing representation error, contact failure, latency miss, safety violation, morphology loss, exception severity, and overhead. The weights are deterministic design choices. They are not claimed to be universal; they make the reporting discipline explicit.

\section{Ablation: Removing Contact Coverage}

If contact coverage were removed from the router, the expected pattern is that contact failure would rise toward the scale-only or guarded-policy level while latency and safety behavior would remain relatively strong. This ablation would test whether the router's advantage is truly multi-axis.

The current table already gives a partial version of this ablation. The latency/safety guarded policy has low safety violation but worse contact failure than the contact-instrumented policy. The contact-instrumented policy has low contact failure but worse safety violation. The router combines both kinds of coverage.

\section{Ablation: Removing Safety Coverage}

Removing safety coverage should reduce utility most in safety-constrained and human-shared workspace conditions. Success might remain high, which is exactly why safety must be reported separately.

This ablation is important for robotics submissions because reviewers often ask whether safety is just another failure. The benchmark says no. Safety changes the objective, and a policy can look good under success while failing under utility.

\section{Ablation: Removing Latency Coverage}

Removing latency coverage should hurt latency-sensitive and mobile manipulation conditions. The failure mode is not a wrong semantic prediction. It is acting on a stale state.

This ablation would be especially useful for real robot logs. The log should record observation timestamp, action timestamp, contact timestamp, and controller update rate. Without these fields, a latency exception can be misclassified as a perception failure.

\section{Ablation: Removing Morphology Coverage}

Removing morphology coverage should hurt locomotion, dexterous transfer, and tool-use conditions. Embodiment-diverse scaling should improve this axis, but it should not automatically recover contact or latency errors.

This ablation distinguishes cross-task generalization from cross-body generalization. A policy can generalize across tasks on one body and still fail when the control regime changes.

\section{Ablation: Removing Router Selection}

If the router were replaced by a fixed average intervention, the expected result would be lower utility. Some conditions need contact instrumentation, others need latency compensation, and others need safety guarding. A fixed intervention pays overhead in easy regimes and misses the active bottleneck in hard regimes.

This is why the paper emphasizes routing rather than simply adding every physical module to every policy. The contribution is not more machinery everywhere; it is condition-aware physical bottleneck selection.

\section{Alternative Metric: Regret To Oracle}

Another way to read the benchmark is regret to the oracle bottleneck router. Scale-only has large regret because its physical residuals are high. The regime-aware router has lower regret but still leaves headroom. That headroom can be assigned to imperfect bottleneck inference, incomplete instrumentation, and intervention overhead.

Regret-to-oracle is useful for future work because it turns the oracle from an unrealistic policy into a diagnostic target. A learned router can be evaluated by how much oracle regret it closes.

\section{Alternative Metric: Physical Residual}

Physical residual is the sum of contact failure, latency miss, safety violation, and morphology loss. This metric abstracts away representation error and asks how much physical error remains after scaling.

The full paper reports the components separately because the components lead to different interventions. However, a physical residual score can be useful for quick dashboards or large benchmark leaderboards.

\section{Alternative Metric: Conditional Slope}

A conditional slope would fit model-scale improvement within each physical regime. A representation-smooth slope should be steeper than a mixed-shift slope. The benchmark does not fit explicit exponents because the scale levels are abstract, but the model-scale table shows the qualitative pattern: success rises slowly while saturation rises sharply.

Future empirical work should report slopes separately for contact, latency, safety, and morphology slices. A single global exponent is too coarse for embodied systems.

\section{Logging Schema For Real Robots}

A real robot version should log task ID, policy ID, model scale, data scale, embodiment source, physical regime label, contact state, force or tactile summary, observation timestamp, action timestamp, safety margin, morphology descriptor, success, and failure reason.

The failure reason should not be a single text field. It should include structured flags for representation error, contact failure, latency miss, safety violation, morphology mismatch, and unknown. This structure is what allows regime-conditioned scaling analysis.

\section{Physical Regime Annotation}

Physical regime labels can be assigned manually at first. A human annotator can tag episodes as smooth, contact-rich, latency-sensitive, safety-constrained, morphology-shifted, or mixed. Later work can learn the regime label.

The important point is that labels must be attached before aggregation. If all episodes are merged into one score, the scaling curve may look smooth while hiding the exception boundary.

\section{Contact Annotation}

Contact annotation should include whether contact was expected, whether contact mode changed, whether force exceeded a threshold, whether slip occurred, and whether tactile or force evidence was available to the policy.

This is not excessive bookkeeping. Without it, a failed insertion can be mislabeled as a perception or planning failure when the actual bottleneck was contact state.

\section{Latency Annotation}

Latency annotation should include sensor delay, policy inference delay, actuator delay, and control update interval. The key derived value is the age of the observation used for the action.

For long-horizon or mobile tasks, this age can dominate performance. A scale law that ignores observation age can overestimate the benefit of model size.

\section{Safety Annotation}

Safety annotation should include minimum margin to human or obstacle, force margin, speed margin, and whether the policy entered a guarded mode. A task can be successful and still unsafe if the minimum margin is violated.

This is why safety violation appears as its own metric. Safety is not just lower success; it is a different constraint on acceptable action.

\section{Morphology Annotation}

Morphology annotation should include body type, kinematic structure, actuator limits, sensor suite, end effector, and whether the training embodiments include a matched control regime.

A broad embodiment dataset should be described in terms of regime coverage, not only number of bodies. Ten bodies in the same control regime may be less diverse than three bodies in distinct regimes.

\section{Protocol: Smooth Regime Study}

The smooth-regime study should use tasks with low contact, low latency, loose safety margins, and matched morphology. This is where scaling should look strongest. If scale does not help here, the policy or dataset may be underpowered.

This slice is the control condition. It prevents the paper from defining exceptions so broadly that scale has no chance to succeed.

\section{Protocol: Contact Regime Study}

The contact-regime study should use insertion, grasp stabilization, or tool alignment with explicit contact mode changes. It should compare a scale-only policy to a contact-instrumented policy and a router.

The central metric should be contact failure, not only task success. A visual policy may get close to the right action but fail when friction or contact mode changes.

\section{Protocol: Latency Regime Study}

The latency-regime study should use a task where state changes during inference or actuation. Mobile manipulation, human-shared workspaces, and dynamic object interaction are natural candidates.

The study should sweep artificial delay as well as real system delay. This makes the conditional curve visible and lets the experimenter distinguish representation failure from stale-state failure.

\section{Protocol: Safety Regime Study}

The safety-regime study should hold task difficulty roughly constant while tightening safety margins. The policy may still reach the goal, but utility should fall if violation risk rises.

This study is important because safety constraints often disappear in aggregate benchmark scores. The paper's reporting rule forces them back into the curve.

\section{Protocol: Morphology Regime Study}

The morphology-regime study should train or evaluate across bodies with known control differences. It should include matched and mismatched morphology conditions, not only a large mixed dataset.

The key question is whether embodiment diversity covers the actual target regime. If it does, diversity should help. If it does not, more bodies may not fix the failure.

\section{Protocol: Mixed Regime Study}

The mixed-regime study should activate more than one bottleneck at a time. For example, a human-shared manipulation task can combine contact, latency, and safety pressure. A locomotion transfer task can combine morphology and latency.

Mixed regimes are where aggregate scale-only curves are most likely to mislead. The benchmark shows scale-only utility near zero in this slice.

\section{Reporting Template}

A paper reporting embodied scaling should include a table with rows for physical regime and columns for model scale, data scale, success, safety violation, contact failure, latency miss, morphology loss, and utility. If a paper cannot fill all columns, it should state which physical axes are unmeasured.

This template makes the hidden assumptions visible. It also helps compare papers that use different robots or policy classes.

\section{Leaderboard Implication}

A leaderboard can encourage overaggregation. If it reports one score, methods optimize for the average regime. A regime-conditioned leaderboard would report separate smooth, contact, latency, safety, morphology, and mixed scores.

The best method might differ by regime. That is not a problem. It is the information the field needs.

\section{Training Implication}

Training should not only scale data quantity. It should scale coverage of physical regimes. A dataset with more contact transitions, more latency variation, more safety-critical examples, and more morphology diversity may matter more than a larger smooth dataset.

This implication is positive: it gives scaling work a sharper target. The right question is what kind of data is being scaled.

\section{Architecture Implication}

Architectures should expose physical bottleneck state to the action policy. This can be done through tactile encoders, force summaries, latency-aware state estimators, safety shields, morphology embeddings, or routers.

The benchmark does not prescribe one architecture. It says that the architecture should make the active bottleneck visible to action selection.

\section{Evaluation Implication}

Evaluation should report mechanism metrics even when success is high. A high success policy with high safety violation is not ready. A high success policy with high contact failure may be relying on easy contact conditions. A high success policy with high morphology loss may not transfer.

Mechanism metrics make scaling curves more honest and more useful.

\section{Why The Corpus Still Matters}

The corpus sweep remains useful as a map of field attention. It suggests that robotics papers already discuss contact, morphology, safety, and physical interaction frequently. It does not prove the benchmark result, but it explains why the benchmark axes are not arbitrary.

V2 sensitivity is also useful because it kept the claim honest. Without that audit, the paper might have overclaimed from keyword ratios.

\section{Why The Benchmark Is Large}

The benchmark is large because exception claims are easy to overfit. A small set of handpicked examples could show scale failure by construction. The full grid crosses several tasks, regimes, scales, data levels, diversity levels, contact levels, latency levels, safety levels, and policies.

The large grid does not make the synthetic assumptions true. It makes the consequences of those assumptions harder to hide.

\section{Why The Benchmark Is Deterministic}

Determinism makes the artifact easier to audit. Every number is a function of explicit factors and formulas. A reviewer can inspect the code and rerun the suite. There is no hidden stochastic lucky run.

For a future hardware benchmark, determinism would be impossible. But the same spirit applies: log enough structure that aggregate results can be audited.

\section{Known Failure Of The Benchmark}

The benchmark may overstate router strength because the router receives clean factor information. In reality, a robot may not know whether the current failure is contact, latency, safety, morphology, or a mixture. Learned bottleneck inference is itself a hard problem.

This limitation is why the oracle and router are separated. The router is not the final answer; it is a target pattern for future systems.

\section{How To Improve The Benchmark}

The next benchmark should include learned regime inference, noisy physical sensors, correlated factor distributions, and real robot traces. It should also include energy cost, wear, calibration drift, and human preference.

These improvements would make the benchmark harder and more realistic. They would not change the main reporting rule: condition scaling curves on physical regime.

\section{One-Sentence Claim}

The one-sentence claim is: model and data scale help embodied policies, but their curves should be reported per physical regime because contact, latency, safety, and morphology can become the dominant bottlenecks.

This sentence is intentionally narrow. It leaves room for strong scaling results while preventing aggregate overclaim.

\section{Final Reviewer Checklist}

A reviewer should ask four questions. Does the paper show where scale helps? Yes: smooth regimes and the model-scale curve. Does the paper show where scale saturates? Yes: mixed physical shift and rising saturation. Does it provide a stronger alternative? Yes: the regime-aware router. Does it overclaim hardware readiness? No: the limitations and guardrails are explicit.
